\chapter{}

Aveva già perduto un anno di tempo, alle belle arti, senza che ancora fosse deciso sul suo conto; il che doveva dipendere dai diversi pareri dei più vecchi avventori, e da suo padre che se ne ricordava soltanto molto di rado e con rabbia. Anna insisteva con pazienza, anche dopo l'infelice prova del disegno, persuasa ch'egli fosse intelligente. Ma era destino che non potesse in alcun modo fargli del bene.

Una mattina decise di portarlo dal parroco, perchè la consigliasse. Aveva già preparato il suo più bel vestito, e voleva far lesta perchè il marito non lo risapesse: ci andava quasi di nascosto. All'improvviso, sentì chiudersi il cuore sempre più stretto; ma non poteva gridare. Non s'accorse nè meno di cadere.

Fu trovata con la testa sul pavimento, verso l'armadio che aveva aperto; tutta stesa in avanti; come quegli animali che hanno avuto una calcagnata sul capo; con gli occhi mezzo schiusi e pieni ancora di vita, con il viso un poco contratto, quasi che le rincrescesse della sua morte soltanto per gli altri, chiedendo di non esserne rimproverata; con una preoccupazione indescrivibile e dolorosa.

Rebecca, ch'era andata a cercarla per ravviarle i capelli, fu la prima a vederla. Ella aprì subito le boccette che servivano quando si trattava delle convulsioni, ma Anna non respirava più.

— Signora padrona! Padrona!

Spaventata e tremando tutta, corse in cucina e s'affacciò a gridare dalla finestra che rispondeva dinanzi all'uscio della trattoria. La intese un cameriere:

— Il padrone! Che venga subito!

Il cameriere, credendo che fosse un attacco di convulsioni più forte del consueto, posò il cencio che aveva in mano e andò in cucina:

— Dov'è il padrone?

— Non è ancora tornato: è restato a pagare il conto dal droghiere.

— Correte subito a cercarlo! La padrona si sente male!

Lo sguattero, che aveva risposto, posò il coltello con il quale puliva il pesce ammonticchiato dentro l'acquaio e tolto allora allora dalla sporta, si asciugò le mani, ravvolse il grembiule su al legacciolo; ed uscì. Ma non potè trovare subito Domenico, che era andato a fare altre spese.

Quando lo vide, tornarono ambedue quasi correndo. Per le scale, Domenico sbattè contro il medico, suo amico e avventore, che scendeva ad aspettarlo:

— Caro Domenico.... Ascoltate un momento!

Il trattore lo prese per le spalle. Il medico gli allontanò le mani, fermandogli i polsi.

— Domenico, questa volta... Quella povera donna!

Egli gridò:

— Mi lasci! È una convulsione.

Ma si sentì gelare tutto, con un gelo che gli veniva a ondate, dalla cima delle dita e si fermava nel mezzo del capo. Credette, lì per lì, che si trattasse di un turbamento della sua intelligenza; ma il respiro affannoso, a lui che respirava così bene, gli ricordò che la cosa quasi presentita era ormai venuta. Come affrontarla? Come vedere Anna morta? Doveva proprio andarci lui?

E quando entrò nella camera, i muri e le porte traballavano e si spalancavano da sé, credette di non vedere niente. Poi toccò il volto già freddo e un po' rigido; e allora chiuse gli occhi, si buttò sopra la moglie e cominciò a piangere.

I suoi gridi stessi lo facevano tremare.

A poco a poco sentì il suo dolore. Tutta la sua enorme violenza, ora, gli pareva cambiata in paura; gli pareva che Poggio a' Meli fosse trascinato via lontano ed egli non aveva il tempo di far qualche cosa; gli pareva che gli usci della sua trattoria si chiudessero da sè e non volessero esser riaperti; e che Anna avesse tanto sofferto per non poter parlare; e tutto crollava in lui.

Il suo dolore era così pieno che tutti avrebbero dovuto consolarlo! Ora si pentiva di non averle voluto bene abbastanza!

Anna s'era raffreddata a poco a poco; e, avendole qualcuno stese le palpebre, parve insolitamente estranea per la prima volta a tutta la gente che le era attorno.

Qualcuno la prese sotto il mento, e la compianse:

— Chi sa che avrebbe voluto dire! Che passione! Povera donna! Così buona!

Pietro la vide già portata sul letto, senza sapere quel che ne dovesse pensare. Domenico gli parlò soltanto quando qualcuno glielo rammentò. Ma senza nessun affetto; quasi con il bisogno di sfuggirlo. E proprio in quel momento, sperò ancora di più di tenerlo con sè per la trattoria. Continuava intanto a gridare che l'udivano anche dalla strada.

— Sembra che stia per scendere dal letto!

Disse Rebecca.

A un tratto Domenico le si accostò un'altra volta, la toccò su i capelli, fece un gesto di disperazione; ed urlò più forte. Pietro, senza provar niente, all'infuori di una vaga inquietudine, si appoggiò ai guanciali e cercò di piangere: dentro di sè chiedevasi se anche gli altri sentissero così poco e provò una consolazione indefinibile quando il padre fu allontanato in modo ch'egli non vide e non udì più il suo dolore; che gli era antipatico come le sue collere.

Rebecca gli disse:

— Povera mamma, voleva tanto bene a te!

A lui gliene importava poco, anzi s'ebbe a male di queste parole; e si allontanò per distrarsi, vergognandosi.

La mattina dell'esequie s'era dimenticato di tutto, quando intravide dall'uscio mezzo aperto il padre che gli si avvicinava. Ebbe, senza spiegarsi il perchè, paura d'esser percosso a sangue.

Domenico gli disse:

— Vestiti; tra poco porteranno via la tua povera mamma.

Pietro si sforzò d'obbedire. Piuttosto, era ora spaventato di qualche sciagura che dovesse capitare a lui!

Discese dal letto; e, fingendo a sé stesso, si vestì cercando d'imitare i gesti di dolore che aveva veduti.

In tal modo finì con il sentire una ilarità muta, mista a terrore.

Ma, quando gli fecero baciare la mamma, prima che la mettessero dentro la cassa, pensò: «Perchè non c'entro anch'io? Metteteci me».

Poi l'assalì uno sgomento inaudito. «Credete che sia morta? Fingete tutti. Anche questa è una finzione. Lo sapevo che m'avreste dato qualche dispiacere violento; e non lo merito.»

Singhiozzò, invaso da una cupa disperazione. Perchè non gli avevano detto prima ch'era morta?

Restò tra le persone che mettevano il cadavere dentro la cassa; ma non avrebbe toccato né meno il lembo della veste. E si meravigliò che gli altri facessero tutto come se si trattasse di una faccenda qualsiasi, con le lacrime e con quei segni di affetto che non sembravano mai finiti: raddrizzare la testa sopra il cuscino scelto con le cifre ricamate, accostare i piedi insieme, accomodare sui capelli un fiore scivolato tra unaspalla e la cassa.

Egli avrebbe voluto che nessuno fosse stato lì; e gli facevano male tutte quelle mani, che si muovevano in fretta. Quelle mani, quelle mani!

Voleva gridare: «Portatela via presto! Perchè non l'avete portata via? Non ce la voglio più in casa». E si meravigliò del padre, che non s'impazientiva, un poco calmato da tutte quelle attenzioni.

Volle seguire il trasporto al cimitero in carrozza chiusa, tirando giù nervosamente le vecchie tendine di seta turchina per non esser visto da nessuno; mentre Domenico anche per risparmio avrebbe voluto andare a piedi. Ma Pietro si preoccupava della gente ferma a guardare nella strada e perfino dinanzi all'uscio di casa. Notò che si alzavano in piedi ed allungavano il collo per veder meglio.
